% Source-faithful mathematical blueprint of Hartshorne, Algebraic Geometry (1977), Appendix C.
% Authorial exposition, examples, and exercises are omitted; mathematical definitions, statements,
% formulas, and proof arguments are retained.
\section{The Zeta Function and the Weil Conjectures}
\label{ha-appC-sec1}
\source{hartshorne-algebraic-geometry:C.1}

\begin{definition}[Zeta function of a variety over a finite field]
\source{hartshorne-algebraic-geometry:C.1}
Let $k=\mathbb F_q$, let $X$ be of finite type over $k$, and put
$N_r=\#X(\mathbb F_{q^r})$.  Define
\[
Z(X,t)=\exp\left(\sum_{r\ge1}N_r\frac{t^r}{r}\right)\in\mathbb Q[[t]].
\]
\end{definition}

\begin{conjecture}[Weil rationality]
\label{ha-appC-conj1-1}
\source{hartshorne-algebraic-geometry:C.1.1}
\dcref{C.1.1}
For smooth projective $X/k$ of dimension $n$, $Z(X,t)$ is rational.
\end{conjecture}

\begin{conjecture}[Weil functional equation]
\label{ha-appC-conj1-2}
\source{hartshorne-algebraic-geometry:C.1.2}
\dcref{C.1.2}
If $E$ is the self-intersection number of a hyperplane section of $X$, then
$Z(1/(q^nt))=\pm q^{nE/2}t^E Z(t)$.
\end{conjecture}

\begin{conjecture}[Weil weights]
\label{ha-appC-conj1-3}
\source{hartshorne-algebraic-geometry:C.1.3}
\dcref{C.1.3}
There are integer polynomials $P_i(t)$, with $P_0(t)=1-t$ and
$P_{2n}(t)=1-q^nt$, such that
$Z(t)=P_1(t)P_3(t)\cdots P_{2n-1}(t)/(P_0(t)P_2(t)\cdots P_{2n}(t))$,
and every reciprocal root of $P_i$ is an algebraic integer of absolute value
$q^{i/2}$.
\end{conjecture}

\begin{conjecture}[Betti numbers]
\label{ha-appC-conj1-4}
\source{hartshorne-algebraic-geometry:C.1.4}
\dcref{C.1.4}
Assuming C.1.3, $B_i=\deg P_i$ are the Betti numbers and
$E=\sum_i(-1)^iB_i$; these agree with the topological Betti numbers when
$X$ is lifted to characteristic zero.
\end{conjecture}

\section{History of Work on the Weil Conjectures}
\label{ha-appC-sec2}
\source{hartshorne-algebraic-geometry:C.2}

The curve case was proved by Weil using correspondences and Jacobians; the
higher-dimensional rationality and functional equation were proved by Dwork.
The search for a characteristic-$p$ cohomology with characteristic-zero
coefficients led to étale and $\ell$-adic cohomology, and Deligne proved the
weight assertion in 1973.

\section{The $\ell$-adic Cohomology}
\label{ha-appC-sec3}
\source{hartshorne-algebraic-geometry:C.3}

\begin{definition}[Etale cohomology and Frobenius data]
\source{hartshorne-algebraic-geometry:C.3}
Fix $\ell\ne\operatorname{char}k$.  For $X$ over $k$ write
\[
H^i(X,\mathbb Q_\ell)=
\left(\varprojlim_m H^i_{\acute et}(X,\mathbb Z/\ell^m)\right)\otimes_{\mathbb Z_\ell}\mathbb Q_\ell.
\]
For a smooth projective $X/\mathbb F_q$, let $F:X\to X$ be Frobenius; its
fixed points satisfy $\operatorname{Fix}(F^r)=X(\mathbb F_{q^r})$.
\end{definition}
\begin{convention}[Finite-dimensionality]
\label{ha-appC-prop3-1}
\source{hartshorne-algebraic-geometry:C.3.1}
\dcref{C.3.1}
$H^i(X,\mathbb Q_\ell)$ is finite-dimensional and vanishes for $i<0$ and
$i>2n$ when $\dim X=n$.
\end{convention}

\begin{convention}[Contravariance]
\label{ha-appC-prop3-2}
\source{hartshorne-algebraic-geometry:C.3.2}
\dcref{C.3.2}
$H^i(-,\mathbb Q_\ell)$ is contravariant in $X$.
\end{convention}

\begin{convention}[Cup product]
\label{ha-appC-prop3-3}
\source{hartshorne-algebraic-geometry:C.3.3}
\dcref{C.3.3}
There is a cup product
$H^i(X,\mathbb Q_\ell)\otimes H^j(X,\mathbb Q_\ell)\to H^{i+j}(X,\mathbb Q_\ell)$.
\end{convention}

\begin{convention}[Poincare duality]
\label{ha-appC-prop3-4}
\source{hartshorne-algebraic-geometry:C.3.4}
\dcref{C.3.4}
If $X$ is smooth proper of dimension $n$, $H^{2n}(X,\mathbb Q_\ell)$ is
one-dimensional and cup product gives a perfect pairing
$H^i(X)\times H^{2n-i}(X)\to H^{2n}(X)$.
\end{convention}

\begin{convention}[Lefschetz fixed point formula]
\label{ha-appC-prop3-5}
\source{hartshorne-algebraic-geometry:C.3.5}
\dcref{C.3.5}
If $f:X\to X$ has isolated multiplicity-one fixed points, then
$L(f,X)=\sum_i(-1)^i\operatorname{Tr}(f^*;H^i(X,\mathbb Q_\ell))$.
\end{convention}

\begin{convention}[Smooth proper base change]
\label{ha-appC-prop3-6}
\source{hartshorne-algebraic-geometry:C.3.6}
\dcref{C.3.6}
For a smooth proper morphism with connected base, $\dim H^i(X_y,\mathbb Q_\ell)$
is constant; in particular it is invariant under finite-field extension.
\end{convention}

\begin{convention}[Comparison over the complex numbers]
\label{ha-appC-prop3-7}
\source{hartshorne-algebraic-geometry:C.3.7}
\dcref{C.3.7}
For smooth proper $X/\CC$,
$H^i(X,\mathbb Q_\ell)\otimes_{\mathbb Q_\ell}\CC\simeq H^i(X_h,\CC)$.
\end{convention}

\begin{convention}[Cycle classes]
\label{ha-appC-prop3-8}
\source{hartshorne-algebraic-geometry:C.3.8}
\dcref{C.3.8}
A codimension-$q$ cycle has a class $\eta(Z)\in H^{2q}(X,\mathbb Q_\ell)$;
rational equivalence is killed, intersections map to cup products, and the
class of a closed point is nonzero.
\end{convention}

\section{Cohomological Interpretation of the Weil Conjectures}
\label{ha-appC-sec4}
\source{hartshorne-algebraic-geometry:C.4}

Let $X/\mathbb F_q$ be smooth projective and let $F:X\to X$ be Frobenius.
The fixed points of $F^r$ are exactly $X(\mathbb F_{q^r})$, so by 3.5,
\[
N_r=\sum_{i=0}^{2n}(-1)^i\operatorname{Tr}(F^{r*};H^i(X,\mathbb Q_\ell)).
\]

\begin{lemma}[Determinant identity]
\label{ha-appC-lem4-1}
\dcref{C.4.1}
For an endomorphism $\varphi$ of a finite-dimensional vector space $V$,
\[
\exp\left(\sum_{r\ge1}\operatorname{Tr}(\varphi^r;V)\frac{t^r}{r}\right)
=\det(1-t\varphi;V)^{-1}.
\]
\source{hartshorne-algebraic-geometry:C.4.1}
\end{lemma}

\begin{proof}
For $\dim V=1$ this is the geometric-series identity.  Induct on dimension
using an invariant subspace and the resulting short exact sequence; traces and
determinants are additive and multiplicative, respectively.
\end{proof}



\begin{theorem}[Cohomological zeta formula]
\label{ha-appC-thm4-2}
\dcref{C.4.2}
Put $P_i(t)=\det(1-tF^*;H^i(X,\mathbb Q_\ell))$.  Then
\[
Z(X,t)=\prod_{i=0}^{2n}P_i(t)^{(-1)^{i+1}}.
\]
Consequently $Z(X,t)$ is rational.  Moreover $P_0(t)=1-t$ and
$P_{2n}(t)=1-q^nt$.
\source{hartshorne-algebraic-geometry:C.4.2}
\uses{ha-appC-lem4-1}
\end{theorem}
\begin{proof}
Insert the Lefschetz trace expression for $N_r$ into the definition of $Z$ and
apply Lemma 4.1 to each cohomological degree.  The Frobenius acts as the
identity on $H^0$, while it acts as multiplication by $q^n$ on the one-
dimensional top cohomology, giving the two displayed endpoint factors.
\end{proof}

\begin{lemma}[Perfect-pairing determinant identity]
\label{ha-appC-lem4-3}
\dcref{C.4.3}
If $V\times W\to K$ is a perfect pairing and endomorphisms $\varphi,\psi$
satisfy $\langle\varphi v,\psi w\rangle=\lambda\langle v,w\rangle$,
then
\[
\det(1-t\psi;W)=(-1)^rt^r\lambda^r
\frac{\det(1-(\lambda t)^{-1}\varphi;V)}{\det(\varphi;V)}.
\]
\source{hartshorne-algebraic-geometry:C.4.3}

\end{lemma}
\begin{proof}
Choose dual bases for the perfect pairing.  The relation
$\langle\varphi v,\psi w\rangle=\lambda\langle v,w\rangle$ says that the matrix of $\psi$ is
$\lambda$ times the inverse transpose of the matrix of $\varphi$.  Taking determinants of
$1-t\psi$ and clearing powers of $t$ gives the displayed identity.
\end{proof}

\begin{theorem}[Functional equation]
\label{ha-appC-thm4-4}
\dcref{C.4.4}
Under the hypotheses of Theorem 4.2,
\[
Z(1/(q^nt))=\pm q^{nE/2}t^E Z(t),
\qquad E=\sum_i(-1)^i\dim H^i(X,\mathbb Q_\ell).
\]
\source{hartshorne-algebraic-geometry:C.4.4}
\uses{ha-appC-thm4-2,ha-appC-lem4-3}
\end{theorem}

\begin{proof}
Apply Poincaré duality to pair $H^i$ with $H^{2n-i}$ and use Lemma 4.3
to relate the corresponding characteristic polynomials.  Substitution in the
product of Theorem 4.2 yields the stated equation, with the sign determined by
the determinants.
\end{proof}



\begin{theorem}[Deligne's theorem]
\label{ha-appC-thm4-5}
\dcref{C.4.5}
The reciprocal roots of each $P_i(t)$ are algebraic integers $\alpha_{ij}$
with $|\alpha_{ij}|=q^{i/2}$.  Hence the Weil conjectures all hold.
\source{hartshorne-algebraic-geometry:C.4.5}
\end{theorem}
\begin{proof}
This is Deligne's theorem, cited in the source; the source explicitly does not
give its proof.  It supplies the algebraic-integrality and absolute-value
assertions for the Frobenius reciprocal roots.
\end{proof}
